\chapter*{Danh mục từ viết tắt và thuật ngữ}
\addcontentsline{toc}{chapter}{Danh mục từ viết tắt và thuật ngữ}
\markboth{Danh mục từ viết tắt}{}

\begin{xltabular}{\linewidth}{|L{3.2cm}|Y|}
  \hline
  \headrow \hd{Viết tắt / thuật ngữ} & \hd{Ý nghĩa} \\ \hline
  \endhead
  API & Application Programming Interface – giao diện lập trình ứng dụng (ở đây là REST/JSON) \\ \hline
  BR & Business Rule – quy tắc nghiệp vụ, đánh mã BR-01 … BR-08 \\ \hline
  CEFR & Common European Framework of Reference – khung tham chiếu năng lực ngôn ngữ châu Âu (A1 … C2) \\ \hline
  CSDL & Cơ sở dữ liệu \\ \hline
  ERD & Entity Relationship Diagram – biểu đồ thực thể liên kết \\ \hline
  FR / NFR & Functional / Non-functional Requirement – yêu cầu chức năng / phi chức năng \\ \hline
  GV, GVHD & Giảng viên, giảng viên hướng dẫn \\ \hline
  HV & Học viên \\ \hline
  KG & Khai giảng \\ \hline
  OOAD & Object-Oriented Analysis and Design – phân tích thiết kế hướng đối tượng \\ \hline
  Placement test & Bài kiểm tra trình độ đầu vào \\ \hline
  RBAC & Role-Based Access Control – phân quyền theo vai trò \\ \hline
  Tờ trình & Đề nghị điện tử gửi Lãnh đạo khi lớp Offline vượt 30 học viên \\ \hline
  UC & Use Case – ca sử dụng \\ \hline
  UML & Unified Modeling Language – ngôn ngữ mô hình hoá thống nhất \\ \hline
  VSTEP & Vietnamese Standardized Test of English Proficiency – bài thi đánh giá năng lực tiếng Anh theo Khung 6 bậc dùng cho Việt Nam \\ \hline
  WebSocket & Giao thức truyền thông hai chiều thời gian thực giữa trình duyệt và máy chủ \\ \hline
\end{xltabular}
